% Options for packages loaded elsewhere
\PassOptionsToPackage{unicode}{hyperref}
\PassOptionsToPackage{hyphens}{url}
\documentclass[
]{article}
\usepackage{xcolor}
\usepackage[margin=1in]{geometry}
\usepackage{amsmath,amssymb}
\setcounter{secnumdepth}{-\maxdimen} % remove section numbering
\usepackage{iftex}
\ifPDFTeX
  \usepackage[T1]{fontenc}
  \usepackage[utf8]{inputenc}
  \usepackage{textcomp} % provide euro and other symbols
\else % if luatex or xetex
  \usepackage{unicode-math} % this also loads fontspec
  \defaultfontfeatures{Scale=MatchLowercase}
  \defaultfontfeatures[\rmfamily]{Ligatures=TeX,Scale=1}
\fi
\usepackage{lmodern}
\ifPDFTeX\else
  % xetex/luatex font selection
\fi
% Use upquote if available, for straight quotes in verbatim environments
\IfFileExists{upquote.sty}{\usepackage{upquote}}{}
\IfFileExists{microtype.sty}{% use microtype if available
  \usepackage[]{microtype}
  \UseMicrotypeSet[protrusion]{basicmath} % disable protrusion for tt fonts
}{}
\makeatletter
\@ifundefined{KOMAClassName}{% if non-KOMA class
  \IfFileExists{parskip.sty}{%
    \usepackage{parskip}
  }{% else
    \setlength{\parindent}{0pt}
    \setlength{\parskip}{6pt plus 2pt minus 1pt}}
}{% if KOMA class
  \KOMAoptions{parskip=half}}
\makeatother
\usepackage{color}
\usepackage{fancyvrb}
\newcommand{\VerbBar}{|}
\newcommand{\VERB}{\Verb[commandchars=\\\{\}]}
\DefineVerbatimEnvironment{Highlighting}{Verbatim}{commandchars=\\\{\}}
% Add ',fontsize=\small' for more characters per line
\usepackage{framed}
\definecolor{shadecolor}{RGB}{248,248,248}
\newenvironment{Shaded}{\begin{snugshade}}{\end{snugshade}}
\newcommand{\AlertTok}[1]{\textcolor[rgb]{0.94,0.16,0.16}{#1}}
\newcommand{\AnnotationTok}[1]{\textcolor[rgb]{0.56,0.35,0.01}{\textbf{\textit{#1}}}}
\newcommand{\AttributeTok}[1]{\textcolor[rgb]{0.13,0.29,0.53}{#1}}
\newcommand{\BaseNTok}[1]{\textcolor[rgb]{0.00,0.00,0.81}{#1}}
\newcommand{\BuiltInTok}[1]{#1}
\newcommand{\CharTok}[1]{\textcolor[rgb]{0.31,0.60,0.02}{#1}}
\newcommand{\CommentTok}[1]{\textcolor[rgb]{0.56,0.35,0.01}{\textit{#1}}}
\newcommand{\CommentVarTok}[1]{\textcolor[rgb]{0.56,0.35,0.01}{\textbf{\textit{#1}}}}
\newcommand{\ConstantTok}[1]{\textcolor[rgb]{0.56,0.35,0.01}{#1}}
\newcommand{\ControlFlowTok}[1]{\textcolor[rgb]{0.13,0.29,0.53}{\textbf{#1}}}
\newcommand{\DataTypeTok}[1]{\textcolor[rgb]{0.13,0.29,0.53}{#1}}
\newcommand{\DecValTok}[1]{\textcolor[rgb]{0.00,0.00,0.81}{#1}}
\newcommand{\DocumentationTok}[1]{\textcolor[rgb]{0.56,0.35,0.01}{\textbf{\textit{#1}}}}
\newcommand{\ErrorTok}[1]{\textcolor[rgb]{0.64,0.00,0.00}{\textbf{#1}}}
\newcommand{\ExtensionTok}[1]{#1}
\newcommand{\FloatTok}[1]{\textcolor[rgb]{0.00,0.00,0.81}{#1}}
\newcommand{\FunctionTok}[1]{\textcolor[rgb]{0.13,0.29,0.53}{\textbf{#1}}}
\newcommand{\ImportTok}[1]{#1}
\newcommand{\InformationTok}[1]{\textcolor[rgb]{0.56,0.35,0.01}{\textbf{\textit{#1}}}}
\newcommand{\KeywordTok}[1]{\textcolor[rgb]{0.13,0.29,0.53}{\textbf{#1}}}
\newcommand{\NormalTok}[1]{#1}
\newcommand{\OperatorTok}[1]{\textcolor[rgb]{0.81,0.36,0.00}{\textbf{#1}}}
\newcommand{\OtherTok}[1]{\textcolor[rgb]{0.56,0.35,0.01}{#1}}
\newcommand{\PreprocessorTok}[1]{\textcolor[rgb]{0.56,0.35,0.01}{\textit{#1}}}
\newcommand{\RegionMarkerTok}[1]{#1}
\newcommand{\SpecialCharTok}[1]{\textcolor[rgb]{0.81,0.36,0.00}{\textbf{#1}}}
\newcommand{\SpecialStringTok}[1]{\textcolor[rgb]{0.31,0.60,0.02}{#1}}
\newcommand{\StringTok}[1]{\textcolor[rgb]{0.31,0.60,0.02}{#1}}
\newcommand{\VariableTok}[1]{\textcolor[rgb]{0.00,0.00,0.00}{#1}}
\newcommand{\VerbatimStringTok}[1]{\textcolor[rgb]{0.31,0.60,0.02}{#1}}
\newcommand{\WarningTok}[1]{\textcolor[rgb]{0.56,0.35,0.01}{\textbf{\textit{#1}}}}
\usepackage{graphicx}
\makeatletter
\newsavebox\pandoc@box
\newcommand*\pandocbounded[1]{% scales image to fit in text height/width
  \sbox\pandoc@box{#1}%
  \Gscale@div\@tempa{\textheight}{\dimexpr\ht\pandoc@box+\dp\pandoc@box\relax}%
  \Gscale@div\@tempb{\linewidth}{\wd\pandoc@box}%
  \ifdim\@tempb\p@<\@tempa\p@\let\@tempa\@tempb\fi% select the smaller of both
  \ifdim\@tempa\p@<\p@\scalebox{\@tempa}{\usebox\pandoc@box}%
  \else\usebox{\pandoc@box}%
  \fi%
}
% Set default figure placement to htbp
\def\fps@figure{htbp}
\makeatother
\setlength{\emergencystretch}{3em} % prevent overfull lines
\providecommand{\tightlist}{%
  \setlength{\itemsep}{0pt}\setlength{\parskip}{0pt}}
\usepackage{booktabs}
\usepackage{caption}
\usepackage{longtable}
\usepackage{colortbl}
\usepackage{array}
\usepackage{anyfontsize}
\usepackage{multirow}
\usepackage{bookmark}
\IfFileExists{xurl.sty}{\usepackage{xurl}}{} % add URL line breaks if available
\urlstyle{same}
\hypersetup{
  pdftitle={STAT332 Statistical Analysis of Networks Sec 01 - Lab 1 - Prof.~Joe Reid},
  pdfauthor={Andrew Sparkes},
  hidelinks,
  pdfcreator={LaTeX via pandoc}}

\title{STAT332 Statistical Analysis of Networks Sec 01 - Lab 1 -
Prof.~Joe Reid}
\author{Andrew Sparkes}
\date{2026-06-01}

\begin{document}
\maketitle

\subsection{1.) Setup}\label{setup}

Welcome to Network Analysis in R. R is a statistical programming
language with a vast history of user developed packages that we are able
to use in order to analyze data. In our case, we will be concerned with
loading a few packages that are critical to our work in network analysis
and then using these packages to analyze networks.

Our integrated developer environment is called R Studio (though, you are
allowed to use Positron if you would rather). I'll demonstrate in class
how to create an R environment for your work in this class. The steps
are rather easy:

1.) Under File select New Project

\begin{figure}
\centering
\pandocbounded{\includegraphics[keepaspectratio,alt={Image showing the new project dialog}]{New_Project_1.png}}
\caption{Image showing the new project dialog}
\end{figure}

2.) Select New Directory or Existing Directory (I would suggest creating
a new directory)

\begin{figure}
\centering
\pandocbounded{\includegraphics[keepaspectratio,alt={Image showing dialog for project type}]{New_Project_2.png}}
\caption{Image showing dialog for project type}
\end{figure}

3.) Select ``New Project''

\begin{figure}
\centering
\pandocbounded{\includegraphics[keepaspectratio,alt={Image showing the ``Create New Project'' Dialog}]{New_Project_3.png}}
\caption{Image showing the ``Create New Project'' Dialog}
\end{figure}

4.) Enter the name of your new directory (You can browse to make it a
subdirectory).

Make sure that you select ``Use renv with this project) which will
create a new R environment for your project. You will want to load up
this environment every time you work on your network data analysis.
Creating a git repo for your work is optional.

When you open up the folder you starte the new project in, you will see
something like this:\\
\strut \\
\pandocbounded{\includegraphics[keepaspectratio,alt={Image showing the folder containing my R project}]{New_Project_4.png}}

Obviously, I have a lot more folders than you are likely to need, but
some of these are auto-created. When you load R-Studio using the
highlighted link, it will open your project which contains all of the
libraries that you will download and use for your work in network
analysis (these are located in the renv file). We will be installing
packages through the package manager:

\begin{figure}
\centering
\pandocbounded{\includegraphics[keepaspectratio,alt={Image showing the package manager in R Studio}]{Package_Manager.png}}
\caption{Image showing the package manager in R Studio}
\end{figure}

The install button on the package manager will give us a dialog box to
type in where we will get our packages for this course.

\begin{figure}
\centering
\pandocbounded{\includegraphics[keepaspectratio,alt={Package Manager Dialog}]{Package_Manager_2.png}}
\caption{Package Manager Dialog}
\end{figure}

Your first task will be to download some of the libraries we are going
to be using. You can start with this list:

\begin{itemize}
\item
  igraph
\item
  igraphdata
\end{itemize}

Once you have these installed, we will get started.

\subsection{2.) Working with packages and Graph
Data}\label{working-with-packages-and-graph-data}

The first thing I'd like you to do is to to create a new .Rmd file. This
is called a markdown file. You can place it wherever you would like
within the folder; however, you need to create this file using
File-\textgreater{} New File -\textgreater{} R Markdown once you have
opened R Studio by the link created for your project. You should get
something that looks like this:

\begin{figure}
\centering
\pandocbounded{\includegraphics[keepaspectratio,alt={Layout of the visual editor within R Studio}]{Rmd_layout.png}}
\caption{Layout of the visual editor within R Studio}
\end{figure}

Notice that I have clicked the ``Visual'' tab so that I can write in a
format that's similar to a typical document editor. The title: and
output: sections are part of something called the YAML. This will give
us the output type for the file which is currently html. In future labs
and homework you will need to know how to create these files, so save
this for your documentation.

\subsubsection{2.a) Getting started.}\label{a-getting-started.}

Using ctrl-alt-i you can create a section where you are going to embed R
code. We're going to start by loading in the igraphdata library and
exploring it a little bit.

\begin{Shaded}
\begin{Highlighting}[]
\FunctionTok{library}\NormalTok{(igraphdata)}
\end{Highlighting}
\end{Shaded}

\begin{Shaded}
\begin{Highlighting}[]
\FunctionTok{data}\NormalTok{(}\AttributeTok{package =} \StringTok{"igraphdata"}\NormalTok{)}
\end{Highlighting}
\end{Shaded}

By running the library(igraphdata) command, it will load the igraphdata
package into memory (RAM). The code-chunk of data(package =
``igraphdata'') opens a viewer that shows the data sets that are within
this package. These are all graph-object data (igraph object) that we
are able to use. Let's start with the most historical (and simplest)
dataset, the Bridges of Koenigsberg data.

\begin{Shaded}
\begin{Highlighting}[]
\FunctionTok{data}\NormalTok{(}\StringTok{"Koenigsberg"}\NormalTok{)}
\NormalTok{g }\OtherTok{\textless{}{-}}\NormalTok{ Koenigsberg}
\FunctionTok{rm}\NormalTok{(Koenigsberg)}
\end{Highlighting}
\end{Shaded}

After running this code cell, you will have done three things. First,
you will have attached the Koenigsberg data. Second, you will create a
copy of this data and called it g (because I'm too lazy to keep typing
Koenigsberg every time). And finally, you will have removed the dataset
Koenigsberg from your active memory. You can now see that the Data
window in R Studio now includes the Koenigsberg data called g.
Technically, this object is of type ``list'' however, the class of this
object is ``igraph''. You can verify this:

\begin{Shaded}
\begin{Highlighting}[]
\FunctionTok{print}\NormalTok{(}\FunctionTok{typeof}\NormalTok{(g))}
\end{Highlighting}
\end{Shaded}

\begin{verbatim}
## [1] "list"
\end{verbatim}

\begin{Shaded}
\begin{Highlighting}[]
\FunctionTok{print}\NormalTok{(}\FunctionTok{class}\NormalTok{(g))}
\end{Highlighting}
\end{Shaded}

\begin{verbatim}
## [1] "igraph"
\end{verbatim}

So, it looks like it's time to learn ``igraph''. We need to load the
library:

\begin{Shaded}
\begin{Highlighting}[]
\FunctionTok{library}\NormalTok{(igraph)}
\end{Highlighting}
\end{Shaded}

\begin{verbatim}
## 
## Attaching package: 'igraph'
\end{verbatim}

\begin{verbatim}
## The following objects are masked from 'package:stats':
## 
##     decompose, spectrum
\end{verbatim}

\begin{verbatim}
## The following object is masked from 'package:base':
## 
##     union
\end{verbatim}

Now, let's look at this ``g'' object.\\

\begin{Shaded}
\begin{Highlighting}[]
\NormalTok{g}
\end{Highlighting}
\end{Shaded}

\begin{verbatim}
## This graph was created by an old(er) igraph version.
## i Call `igraph::upgrade_graph()` on it to use with the current igraph version.
## For now we convert it on the fly...
\end{verbatim}

\begin{verbatim}
## IGRAPH 227bd5e UN-- 4 7 -- The seven bidges of Koenigsberg
## + attr: name (g/c), name (v/c), Euler_letter (v/c), Euler_letter (e/c),
## | name (e/c)
## + edges from 227bd5e (vertex names):
## [1] Altstadt-Loebenicht--Kneiphof          
## [2] Altstadt-Loebenicht--Kneiphof          
## [3] Altstadt-Loebenicht--Lomse             
## [4] Kneiphof           --Lomse             
## [5] Vorstadt-Haberberg --Lomse             
## [6] Kneiphof           --Vorstadt-Haberberg
## [7] Kneiphof           --Vorstadt-Haberberg
\end{verbatim}

It's actually a good thing we started with a simple graph here, because
just typing the name of the graph object will produce both a summary of
the graph AND a list of all of the edges. If we just want the summary,
we use the ``summary'' function.

\begin{Shaded}
\begin{Highlighting}[]
\FunctionTok{summary}\NormalTok{(g)}
\end{Highlighting}
\end{Shaded}

\begin{verbatim}
## IGRAPH 227bd5e UN-- 4 7 -- The seven bidges of Koenigsberg
## + attr: name (g/c), name (v/c), Euler_letter (v/c), Euler_letter (e/c),
## | name (e/c)
\end{verbatim}

Before we try to figure out what this output means, let's just get a
quick plot of the data.

\begin{Shaded}
\begin{Highlighting}[]
\FunctionTok{plot}\NormalTok{(g)}
\end{Highlighting}
\end{Shaded}

\pandocbounded{\includegraphics[keepaspectratio]{Lab1_graph_basics_files/figure-latex/unnamed-chunk-8-1.pdf}}

This graph data already exists, so we didn't have to create it
ourselves, but we could have (and will need to) in the future. For now,
we will continue exploring this data and then move on to data creation
in part 3. At this time, we should probably parse out that summary
function. Let's remind ourselves of the output:

\begin{Shaded}
\begin{Highlighting}[]
\FunctionTok{summary}\NormalTok{(g)}
\end{Highlighting}
\end{Shaded}

\begin{verbatim}
## IGRAPH 227bd5e UN-- 4 7 -- The seven bidges of Koenigsberg
## + attr: name (g/c), name (v/c), Euler_letter (v/c), Euler_letter (e/c),
## | name (e/c)
\end{verbatim}

Ignoring the numbers directly after IGRAPH, You can see that the ouput
is: UN-- 4 7 -- The seven bridges of Koenigsberg. The first four
characters indicate the poperties of the graph. U means unidrected and D
means directed. N tells us that a vertex attribute called ``name''
exists (so that there are names for the verticies and numbers shouldn't
be used for the labels). W would indicate that the graph is weighted.
And B as the fourth character would indicate a vertex attributed called
``type'' (B is historic and represents a bipartite graph). Dashes
indicate the absence of these parameters. The number 4 indicates that
this graph has 4 vertices and the number 7 indicates the number of
edges. The next two dashes are just to separate the title of the graph
from the encoded information. The next thing that we can see is a list
of the node and edge attributes of the graph. Anything with a ``v''
indicates that it's an attribute of the vertex and anything with an
``e'' indicates it's an edge attribute. Let's explore these attributes.

\begin{Shaded}
\begin{Highlighting}[]
\FunctionTok{V}\NormalTok{(g)}\SpecialCharTok{$}\NormalTok{name}
\end{Highlighting}
\end{Shaded}

\begin{verbatim}
## [1] "Altstadt-Loebenicht" "Kneiphof"            "Vorstadt-Haberberg" 
## [4] "Lomse"
\end{verbatim}

\begin{Shaded}
\begin{Highlighting}[]
\FunctionTok{V}\NormalTok{(g)}\SpecialCharTok{$}\NormalTok{Euler\_letter}
\end{Highlighting}
\end{Shaded}

\begin{verbatim}
## [1] "B" "A" "C" "D"
\end{verbatim}

\begin{Shaded}
\begin{Highlighting}[]
\FunctionTok{E}\NormalTok{(g)}\SpecialCharTok{$}\NormalTok{Euler\_letter}
\end{Highlighting}
\end{Shaded}

\begin{verbatim}
## [1] "a" "b" "f" "e" "g" "c" "d"
\end{verbatim}

\begin{Shaded}
\begin{Highlighting}[]
\FunctionTok{E}\NormalTok{(g)}\SpecialCharTok{$}\NormalTok{name}
\end{Highlighting}
\end{Shaded}

\begin{verbatim}
## [1] "Kraemer Bruecke" "Schmiedebruecke" "Holzbruecke"     "Honigbruecke"   
## [5] "Hohe Bruecke"    "Gruene Bruecke"  "Koettelbruecke"
\end{verbatim}

We're going to ignore the Euler letters as they aren't particularly
important to this class (but I would strongly recommend looking up the
history of this problem if you are interested!). V(g)\$name gives us the
names of the communities at each end of the bridges. These represent our
vertex names for this graph. The E(g)\$name gives us the names of the
bridges (``Die Bruecke'' means ``The Bridge'' in German, in case you
were wondering.) There are no other edge properties to worry about here,
so we are good to go.

So, from our notes, we can tell that this is an undirected graph. Create
the adjacency matrix for this by hand before we look at how to produce
it with code. For the record, adjacency matrix storage of graph objects
is very memory inefficient, so I really wouldn't do this with a large
graph.

\begin{Shaded}
\begin{Highlighting}[]
\FunctionTok{as\_adjacency\_matrix}\NormalTok{(g, }\AttributeTok{sparse =} \ConstantTok{FALSE}\NormalTok{)}
\end{Highlighting}
\end{Shaded}

\begin{verbatim}
##                     Altstadt-Loebenicht Kneiphof Vorstadt-Haberberg Lomse
## Altstadt-Loebenicht                   0        2                  0     1
## Kneiphof                              2        0                  2     1
## Vorstadt-Haberberg                    0        2                  0     1
## Lomse                                 1        1                  1     0
\end{verbatim}

Notice that the diagonal is 0? What does this mean? Some of the nodes
have multiple paths between them making this a multigraph. If we want to
get the degree's of the nodes for this graph, we can do it using the
``degree'' function.

\begin{Shaded}
\begin{Highlighting}[]
\FunctionTok{degree}\NormalTok{(g)}
\end{Highlighting}
\end{Shaded}

\begin{verbatim}
## Altstadt-Loebenicht            Kneiphof  Vorstadt-Haberberg               Lomse 
##                   3                   5                   3                   3
\end{verbatim}

If you want to extract the edge list, you can do that with the E()
function.

\begin{Shaded}
\begin{Highlighting}[]
\FunctionTok{E}\NormalTok{(g)}
\end{Highlighting}
\end{Shaded}

\begin{verbatim}
## + 7/7 edges from 227bd5e (vertex names):
## [1] Altstadt-Loebenicht--Kneiphof          
## [2] Altstadt-Loebenicht--Kneiphof          
## [3] Altstadt-Loebenicht--Lomse             
## [4] Kneiphof           --Lomse             
## [5] Vorstadt-Haberberg --Lomse             
## [6] Kneiphof           --Vorstadt-Haberberg
## [7] Kneiphof           --Vorstadt-Haberberg
\end{verbatim}

\subsubsection{2b.) Statistics}\label{b.-statistics}

\paragraph{Size}\label{size}

Remember that the size of a graph is simply the number of vertices. This
can be obtained from the summary function, but to isolate this number,
we can use a function: gorder(). For what it's worth, another name for
the size of a graph is the ``order'' of the graph; hence graph-order.

\begin{Shaded}
\begin{Highlighting}[]
\FunctionTok{gorder}\NormalTok{(g)}
\end{Highlighting}
\end{Shaded}

\begin{verbatim}
## [1] 4
\end{verbatim}

\paragraph{Components}\label{components}

In order to count the number of components on a graph, we can do this
with a function called: count\_components()

\begin{Shaded}
\begin{Highlighting}[]
\FunctionTok{count\_components}\NormalTok{(g)}
\end{Highlighting}
\end{Shaded}

\begin{verbatim}
## [1] 1
\end{verbatim}

\paragraph{Isolates}\label{isolates}

Finding the number of isolates is actually a little bit weird. There is
no built in function for this, but recalling that an isolate is a vertex
with degree zero, we can use the degree function to produce a vector of
the nodes and their degrees. What we need to do is to ``count'' the
number of times where the degree is zero. This is done with the sum
function. Hence:

\begin{Shaded}
\begin{Highlighting}[]
\FunctionTok{sum}\NormalTok{(}\FunctionTok{degree}\NormalTok{(g) }\SpecialCharTok{==} \DecValTok{0}\NormalTok{)}
\end{Highlighting}
\end{Shaded}

\begin{verbatim}
## [1] 0
\end{verbatim}

If we want to look at the proportion of values that are isolates, we can
do this using the gorder() function from earlier (yeah, this is the
reason we needed that.)

Proportion of isolates:

\begin{Shaded}
\begin{Highlighting}[]
\FunctionTok{sum}\NormalTok{(}\FunctionTok{degree}\NormalTok{(g) }\SpecialCharTok{==} \DecValTok{0}\NormalTok{) }\SpecialCharTok{/} \FunctionTok{gorder}\NormalTok{(g)}
\end{Highlighting}
\end{Shaded}

\begin{verbatim}
## [1] 0
\end{verbatim}

\paragraph{Density}\label{density}

Let's use the function to find the density of the graph: edge\_density()

\begin{Shaded}
\begin{Highlighting}[]
\FunctionTok{edge\_density}\NormalTok{(g)}
\end{Highlighting}
\end{Shaded}

\begin{verbatim}
## [1] 1.166667
\end{verbatim}

Oh no! Edge density can't be a number greater than 1! What could have
caused this? If a graph contains loops or is a multigraph, the density
calculation will be wrong. We need to adjust for this by looking at a
graph where only the unique edges are present and loops are accounted
for. The edge\_density function can account for the loops using a bool
edge\_density(graph, loops = FALSE). Changing this to true for loops
will adjust the calculation, but it still doesn't fix the other issue.
You can actually perform both tasks at once using the ``simplify''
function.

\begin{Shaded}
\begin{Highlighting}[]
\FunctionTok{edge\_density}\NormalTok{(}\FunctionTok{simplify}\NormalTok{(g))}
\end{Highlighting}
\end{Shaded}

\begin{verbatim}
## [1] 0.8333333
\end{verbatim}

\paragraph{Diameter}\label{diameter}

\begin{Shaded}
\begin{Highlighting}[]
\FunctionTok{diameter}\NormalTok{(g)}
\end{Highlighting}
\end{Shaded}

\begin{verbatim}
## [1] 2
\end{verbatim}

\paragraph{Average Path Length}\label{average-path-length}

igraph used to have a function for average path length called:
average.path.length(). However, this function is now depricated and has
been replaced by the function: mean\_distance()

\begin{Shaded}
\begin{Highlighting}[]
\FunctionTok{mean\_distance}\NormalTok{(g)}
\end{Highlighting}
\end{Shaded}

\begin{verbatim}
## [1] 1.166667
\end{verbatim}

\subsubsection{2c.) Representation of summary
statistics}\label{c.-representation-of-summary-statistics}

Obviously, this is not very effective at presenting the results. We
should make a summary table of these. I'm going to use the packages
data.table and gt for this (if you don't have it, do you remember how to
get it?)

\begin{Shaded}
\begin{Highlighting}[]
\FunctionTok{library}\NormalTok{(data.table)}
\NormalTok{table }\OtherTok{\textless{}{-}} \FunctionTok{data.table}\NormalTok{(}
  \AttributeTok{ID =} \FunctionTok{c}\NormalTok{(}\StringTok{"Size"}\NormalTok{, }\StringTok{"Components"}\NormalTok{, }\StringTok{"Isolate Proportion"}\NormalTok{, }\StringTok{"Density"}\NormalTok{, }\StringTok{"Diameter"}\NormalTok{, }\StringTok{"Average Path Length"}\NormalTok{),}
  \AttributeTok{Value =} \FunctionTok{c}\NormalTok{(}\FunctionTok{gorder}\NormalTok{(g), }\FunctionTok{count\_components}\NormalTok{(g), }\FunctionTok{sum}\NormalTok{(}\FunctionTok{degree}\NormalTok{(g) }\SpecialCharTok{==} \DecValTok{0}\NormalTok{) }\SpecialCharTok{/} \FunctionTok{gorder}\NormalTok{(g), }\FunctionTok{edge\_density}\NormalTok{(}\FunctionTok{simplify}\NormalTok{(g)), }\FunctionTok{diameter}\NormalTok{(g), }\FunctionTok{mean\_distance}\NormalTok{(g))}
\NormalTok{)}
\end{Highlighting}
\end{Shaded}

\begin{Shaded}
\begin{Highlighting}[]
\FunctionTok{library}\NormalTok{(gt)}
\NormalTok{table }\SpecialCharTok{\%\textgreater{}\%} \FunctionTok{gt}\NormalTok{()}
\end{Highlighting}
\end{Shaded}

\begin{table}[t]
\fontsize{12.0pt}{14.0pt}\selectfont
\begin{tabular*}{\linewidth}{@{\extracolsep{\fill}}lr}
\toprule
ID & Value \\ 
\midrule\addlinespace[2.5pt]
Size & 4.0000000 \\ 
Components & 1.0000000 \\ 
Isolate Proportion & 0.0000000 \\ 
Density & 0.8333333 \\ 
Diameter & 2.0000000 \\ 
Average Path Length & 1.1666667 \\ 
\bottomrule
\end{tabular*}
\end{table}

\subsection{3.) Your turn \#1:}\label{your-turn-1}

A collection of journal articles is provided for your consumption. I
have provided csv files for both the nodes and links as well as code to
read and create a network from this data. Using what you know from the
code above, create a plot of this network and a table that summarizes
the statistics describing this data.

\begin{Shaded}
\begin{Highlighting}[]
\NormalTok{nodes }\OtherTok{\textless{}{-}} \FunctionTok{read.csv}\NormalTok{(}\StringTok{"Citation\_vertices.csv"}\NormalTok{, }\AttributeTok{header =} \ConstantTok{TRUE}\NormalTok{, }\AttributeTok{as.is =} \ConstantTok{TRUE}\NormalTok{)}
\NormalTok{links }\OtherTok{\textless{}{-}} \FunctionTok{read.csv}\NormalTok{(}\StringTok{"Citation\_Edgelist.csv"}\NormalTok{, }\AttributeTok{header =} \ConstantTok{TRUE}\NormalTok{, }\AttributeTok{as.is =} \ConstantTok{TRUE}\NormalTok{)}
\NormalTok{net }\OtherTok{\textless{}{-}} \FunctionTok{graph\_from\_data\_frame}\NormalTok{(}\AttributeTok{d =}\NormalTok{ links, }\AttributeTok{vertices =}\NormalTok{ nodes, }\AttributeTok{directed =} \ConstantTok{TRUE}\NormalTok{)}
\end{Highlighting}
\end{Shaded}

\subsubsection{Plot (remember, this is a directed
graph)}\label{plot-remember-this-is-a-directed-graph}

\begin{Shaded}
\begin{Highlighting}[]
\FunctionTok{plot}\NormalTok{(net)}
\end{Highlighting}
\end{Shaded}

\pandocbounded{\includegraphics[keepaspectratio]{Lab1_graph_basics_files/figure-latex/unnamed-chunk-28-1.pdf}}

\subsubsection{Table of Summary
Information}\label{table-of-summary-information}

\begin{Shaded}
\begin{Highlighting}[]
\NormalTok{make\_table\_happen }\OtherTok{\textless{}{-}} \ControlFlowTok{function}\NormalTok{(graph) \{}
  \FunctionTok{data.table}\NormalTok{(}
    \AttributeTok{ID =} \FunctionTok{c}\NormalTok{(}\StringTok{"Size"}\NormalTok{, }\StringTok{"Components"}\NormalTok{, }\StringTok{"Isolate Proportion"}\NormalTok{, }\StringTok{"Density"}\NormalTok{, }\StringTok{"Diameter"}\NormalTok{, }\StringTok{"Average Path Length"}\NormalTok{),}
    \AttributeTok{Value =} \FunctionTok{c}\NormalTok{(}\FunctionTok{gorder}\NormalTok{(graph), }\FunctionTok{count\_components}\NormalTok{(graph), }\FunctionTok{sum}\NormalTok{(}\FunctionTok{degree}\NormalTok{(graph) }\SpecialCharTok{==} \DecValTok{0}\NormalTok{) }\SpecialCharTok{/} \FunctionTok{gorder}\NormalTok{(graph), }\FunctionTok{edge\_density}\NormalTok{(}\FunctionTok{simplify}\NormalTok{(graph)), }\FunctionTok{diameter}\NormalTok{(graph), }\FunctionTok{mean\_distance}\NormalTok{(graph))}
\NormalTok{  ) }\SpecialCharTok{|\textgreater{}}
    \FunctionTok{gt}\NormalTok{()}
\NormalTok{\}}

\FunctionTok{make\_table\_happen}\NormalTok{(net)}
\end{Highlighting}
\end{Shaded}

\begin{table}[t]
\fontsize{12.0pt}{14.0pt}\selectfont
\begin{tabular*}{\linewidth}{@{\extracolsep{\fill}}lr}
\toprule
ID & Value \\ 
\midrule\addlinespace[2.5pt]
Size & 46.00000000 \\ 
Components & 13.00000000 \\ 
Isolate Proportion & 0.19565217 \\ 
Density & 0.02222222 \\ 
Diameter & 8.00000000 \\ 
Average Path Length & 2.84883721 \\ 
\bottomrule
\end{tabular*}
\end{table}

\subsection{4.) Your turn \#2:}\label{your-turn-2}

We are going to generate a random graph based on the Erdos-Renyi model.
Your task is to produce the Plot and table of summary information for
this network as well.

\begin{Shaded}
\begin{Highlighting}[]
\FunctionTok{set.seed}\NormalTok{(}\DecValTok{123}\NormalTok{)}
\NormalTok{g2 }\OtherTok{\textless{}{-}} \FunctionTok{sample\_gnp}\NormalTok{(}\DecValTok{50}\NormalTok{, }\FloatTok{0.048}\NormalTok{)}
\end{Highlighting}
\end{Shaded}

\subsubsection{Plot}\label{plot}

\begin{Shaded}
\begin{Highlighting}[]
\FunctionTok{plot}\NormalTok{(g2)}
\end{Highlighting}
\end{Shaded}

\pandocbounded{\includegraphics[keepaspectratio]{Lab1_graph_basics_files/figure-latex/unnamed-chunk-31-1.pdf}}

\subsubsection{Table of Summary
Information}\label{table-of-summary-information-1}

\begin{Shaded}
\begin{Highlighting}[]
\FunctionTok{make\_table\_happen}\NormalTok{(g2)}
\end{Highlighting}
\end{Shaded}

\begin{table}[t]
\fontsize{12.0pt}{14.0pt}\selectfont
\begin{tabular*}{\linewidth}{@{\extracolsep{\fill}}lr}
\toprule
ID & Value \\ 
\midrule\addlinespace[2.5pt]
Size & 50.00000000 \\ 
Components & 5.00000000 \\ 
Isolate Proportion & 0.08000000 \\ 
Density & 0.04816327 \\ 
Diameter & 10.00000000 \\ 
Average Path Length & 4.32463768 \\ 
\bottomrule
\end{tabular*}
\end{table}

\subsection{}\label{section}

\subsection{5.) Results}\label{results}

Obviously, we have some work that needs done. Our plots don't look very
good, we aren't representing any attributes of the vertices or edges, we
aren't interpreting the statistical measures, we aren't finding clusters
or cliques, we don't know how resistant these networks are, we don't
know anything about the hubs or critical bridges in the data, etc\ldots;
however, we now have a start.

What you need to do nest is to ``knit'' your document into an html using
the ``Knit'' button near the top of the screen. This will give you an
HTML that you can turn in.

\end{document}
